\section{Functional outputs and absolute observation}\label{sec:absolute}

The positive-coordinate invariant and ground-substitution covariance
were stated in Section~\ref{sec:setting}. They permit matching fixed
heads without replacing a source scalar by an invented value.

\begin{lemma}[Absolute observer]\label{lem:absolute}
Let $T:\rho$ have only ground free names $\mathbf z$.  Its entire raw scalar
graph $v=\can{T(\mathbf z)}_0$ is existentially definable by actual ground-source
observations, addition and constants, with all shared names preserved.
\end{lemma}
\begin{proof}
Close the ground names to literal zero, obtaining $T^*$, and evaluate the
closed scalar $C=\can{T^*}_0$.  Put $m=\lv(\rho)$ and $p=2(m+1)+1$.
For a fresh unused ground $u$, the full vector size multiplier gives
\[
 \can{\lambda u:o.T}_0=p\can T_0+p-1.
\]
The positive profiles agree for $T$ and $T^*$.  With fresh positive ground
inputs $x,y$, form the actual ground sources
\[
 G(\mathbf z,x)=(\lambda Y:\rho.0)((\lambda u:o.T)x),\qquad
 G^*(y)=(\lambda Y:\rho.0)((\lambda u:o.T^*)y).
\]
The inner applications and outer erasers use matching fixed heads on the two
sides.  Twice cancelling their injective fixed-tail collapses yields
\[
 G_0=G^*_0\quad\Longleftrightarrow\quad
 p\can T_0+x_0=pC+y_0.
\]
Add $pv+x_0=pC+y_0$.  These two equations force $v=\can T_0$;
conversely raw witnesses $x_0=pC+1$, $y_0=pv+1$ work.  All raw values are
encoded by naturals plus one.  The anchor $C$ is effectively computed by the
closed evaluator; that cost is charged, with no feasible-size claim.
\end{proof}

\begin{theorem}[Exact decision equivalence]\label{thm:decision-equivalence}
Universal weak/strict canonical comparison on the full all-ground-free-variable
domain, its ground-output restriction, and unsatisfiability in $\Esrc$ are
computably equivalent.  Satisfiability uses answer complementation.
\end{theorem}
\begin{proof}
For a general-output comparison first decide its closed positive-coordinate
tests.  If they pass, a countervaluation is an existential inequality between
the scalar observations defined by Lemma~\ref{lem:absolute}.  The Boolean
compiler handles that sentence.  Conversely it reduces arbitrary source
arithmetic unsatisfiability to ground-output comparison.  The ground-output
domain is a special case of the general domain.  The closed-anchor evaluations
in the general-output direction justify computability, not a uniform
polynomial bound.
\end{proof}

This equivalence transfers a decision theorem when one is available.
For example, the ground-input decision theorem for the permitted controls
of \cite{groundcomparison} decides this actual-source arithmetic at those
controls. The equivalence alone does not supply a decision procedure
for the unchanged standard control.


\subsection{Eliminating common functional unit observations}\label{sec:units}

A unit observation substitutes the number one in every coordinate of a free
functional variable.  It is an observation of the assignment, not a claim
that a closed term has the all-one vector.  We first show how to eliminate
such observations without requiring that inhabitant.

\begin{lemma}[One-name scalar elimination]\label{lem:one-unit}
Let $A:\tau$ be an actual source, $h=\can A$, $m=\lv(\tau)$, and let
$X:\sigma$ be one of its free functional names.  Put
\[
 n=\lv(\sigma)+1,\quad d=2^n\sum_{i=0}^ms_X(h_i),\quad
 P=(\partial^X_0h_0)_0,\quad R=h_0\{X:=\mathbf1\},\quad L=\lambda X.A.
\]
The graph of the desired scalar $R$ has a finite positive-coefficient linear
definition using scalar observations of actual sources without free $X$.
\end{lemma}
\begin{proof}
The original abstraction clause gives $L_0=dP+R$.  If the whole $h_0$ is
$X$-free, the stopping rule gives $P=R+1$, hence
$L_0=(d+1)R+d$.  The positive coefficient uniquely determines $R$.

Suppose $h_0$ is active.  Introduce a fresh unused ground $u$ and put
$V=\lambda u:o.A$, $k=2(m+1)$, $p=k+1$, $h'=\can V$,
$d'=2^n\sum_js_X(h'_j)$ and $L'=\lambda X.V$.
The literal scalar is $h'_0=k(h_0+1)+h_0$, where the multiplication means
the original finite sum of copies, before any collecting.  Its partial is
$pP+4k$: a copy of $h_0+1$ contributes $P+3$, and the $k$ sum nodes
joining the copies and retained $h_0$ contribute $k$ more.  The retained
unit value is $pR+k$.  The scalar size is
$ps_X(h_0)+3k>s_X(h_0)$.

No positive-coordinate size decreases.  For $m\ge1$ the new first field
has two copies of $h_1$ and finite constants, and higher fields have two
copies of their old field.  An entire $X$-free field has size one on both
sides; an active field has at least its original size.  For $m=0$ the new
positive field is one.  Therefore $d'>d$.  With $c=4kd'+k$ we obtain
\[
 L'_0=pd'P+pR+c,\qquad
 pd'L_0+dc=dL'_0+p(d'-d)R.
\]
The coefficient of $R$ is positive.  These two sources have no free $X$,
and every other functional name can still receive its common unit observation.
The intended $R$ satisfies the equation and is its unique solution once the
other observations are fixed.  All coefficients are finite syntactic integers.
\end{proof}

\begin{theorem}[Finite unit-observation elimination]\label{thm:unit-elimination}
For an actual source $A$ and a finite set of all its free functional names,
the absolute scalar after their common unit observation is existentially
definable in $\Esrc$, retaining all original shared ground inputs.
\end{theorem}
\begin{proof}
Apply Lemma~\ref{lem:one-unit} to one name, then recursively to the actual
sources $L,L'$ on the remaining set.  Each branch removes a name; the new
$u$ is ground.  Induction on the number of functional names terminates.
Leaves have only ground free names.  Lemma~\ref{lem:absolute} supplies their
absolute scalar graphs, including leaves of functional output type.

For each fixed ground valuation the intended values supply positive witnesses.
Reverse induction through the positive-coefficient equations proves uniqueness
of every internal scalar.  Encoding raw values as naturals plus one merely
changes fixed additive constants.  The Boolean compiler consumes the resulting
finite system.  No repeated original name has acquired an independent value.
\end{proof}

